\section{Results and Analysis}
\label{sec:results}

\subsection{Pilot: Windows tasks}
\label{sec:pilot}
Table~\ref{tab:pilot} reports the pilot on the \ResPilotTasks{} Windows tasks. All tasks completed and every
verdict was correct (accuracy \ResPilotAccuracy{}): the eight attacks were judged malicious and the three
benign look-alikes benign, so the false-positive and false-negative rates were both zero. Mean quality was
\ResPilotQuality{}, mean efficiency \ResPilotEfficiency{} and mean reward \ResPilotReward{}, with
\ResPilotCalls{} tool calls per task (range \ResPilotMinCalls{}--\ResPilotMaxCalls{}). The whole run took
\ResPilotWall\,s of wall-clock time with four tasks in parallel.

\begin{table}[t]\centering\footnotesize
\caption{Pilot run on the 11 Windows tasks (Claude Opus 5.5, effort \emph{high}, full configuration). Quality is 0 for a wrong verdict, so every verdict here was correct.}
\label{tab:pilot}
\resizebox{\columnwidth}{!}{%
\begin{tabular}{@{}llrrrr@{}}\toprule
Task & Expected & Quality & Effic. & Reward & Calls\\\midrule
\texttt{logon-script-persistence} & malicious & 0.917 & 0.714 & 0.864 & 7\\
\texttt{lsass-dump-comsvcs} & malicious & 1.000 & 0.714 & 0.943 & 7\\
\texttt{service-binpath-hijack} & malicious & 0.922 & 0.600 & 0.849 & 10\\
\texttt{hta-startup-folder} & malicious & 0.922 & 0.500 & 0.830 & 12\\
\texttt{firewall-rule-added} & malicious & 0.887 & 0.667 & 0.828 & 6\\
\texttt{bits-download} & malicious & 0.929 & 0.800 & 0.891 & 5\\
\texttt{run-key-and-beacon} & malicious & 0.908 & 0.625 & 0.840 & 8\\
\texttt{lsass-query-by-svchost} & benign & 1.000 & 1.000 & 1.000 & 3\\
\texttt{python-http-server} & malicious & 0.922 & 0.455 & 0.822 & 11\\
\texttt{w32time-runtime-key} & benign & 1.000 & 0.750 & 0.950 & 4\\
\texttt{gpsvc-service-host} & benign & 0.850 & 0.429 & 0.753 & 7\\
\midrule
Mean & & 0.933 & 0.659 & 0.870 & 7.27\\
\bottomrule\end{tabular}}\end{table}


\paragraph{Where reward is lost (RQ2)} Since $R=Q(0.8+0.2E)$, the gap between the mean reward and its maximum
of 1 decomposes into a quality part, $\overline{1-Q}=0.068$, and an efficiency part,
$\overline{0.2\,Q(1-E)}=0.063$. Efficiency therefore costs almost as much as imperfect artifacts.
Fig.~\ref{fig:pilotcalls} shows why: the agent used 80 tool calls against a target of 49 (the sum of
$n^\star+1$), and 10 of 11 tasks exceeded their target. The largest overruns were the HTA execution (12 calls
against a target of 6), the Python HTTP server (11 against 5) and the service-configuration change (10 against
6), all multi-step chains in which the agent fetched relations one at a time. Benign tasks were cheaper
(4.7 calls on average, efficiency 0.73) than malicious ones (8.3 calls, efficiency 0.63), except the Group
Policy service host (7 calls against a target of 3), where the agent kept collecting evidence after the parent
process alone settled the question.

\paragraph{Quality} Quality was 1.0 on three tasks and at least 0.85 on all. With a correct verdict, the
score leaves only two ways to lose quality: for malicious tasks, missing core entities or naming entities
outside the acceptable set; for benign tasks, citing only some anchors or non-decisive fields. The per-call
traces of the study runs attribute these losses per task.

% Pilot: tool calls per task against the optimal count (pgfplots, reads data/pilot_per_task.csv).
\begin{figure}[t]
\centering
\begin{tikzpicture}
\begin{axis}[
  xbar, bar width=4pt, width=0.62\columnwidth, height=7.2cm,
  xmin=0, xmax=13, xlabel={Tool calls}, xmajorgrids, grid style={gray!25},
  symbolic y coords={gpsvc-service-host,w32time-runtime-key,lsass-query-by-svchost,python-http-server,run-key-and-beacon,bits-download,firewall-rule-added,hta-startup-folder,service-binpath-hijack,lsass-dump-comsvcs,logon-script-persistence},
  ytick=data, yticklabel style={font=\tiny\ttfamily}, tick label style={font=\scriptsize},
  legend style={font=\scriptsize, at={(0.98,0.02)}, anchor=south east, draw=gray!40},
  enlarge y limits=0.06,
]
\addplot[fill=cA, draw=none] table[col sep=comma, x=calls, y=task] {data/pilot_per_task.csv};
\addplot[fill=cB, draw=none] table[col sep=comma, x expr=\thisrow{optimal}+1, y=task] {data/pilot_per_task.csv};
\legend{Used, Optimal + 1}
\end{axis}
\end{tikzpicture}
\caption{Pilot run: tool calls used per task against the efficiency target (optimal + 1, at which efficiency
reaches 1). Every verdict was correct; the gap between the bars is the efficiency loss.}
\label{fig:pilotcalls}
\end{figure}


\subsection{Full study and ablation (RQ1, RQ3, RQ4)}
\label{sec:study}
\ifnum\ResStudyAvailable=1
Table~\ref{tab:ablation} reports the full configuration and the five ablations over \ResStudyTasks{} tasks
with \ResStudyRepeats{} repeats per configuration; Table~\ref{tab:effects} gives paired task-level effects,
Table~\ref{tab:breakdown} a breakdown by operating system and expected verdict, Table~\ref{tab:pertask} the
per-task matrix and Table~\ref{tab:safeguards} what happened at the tool interface. Repeat~1 was rescored from
v1.0 (\S\ref{sec:experiments}); all \ResPairedRepeats{} repeats have per-task results and traces.

\begin{table*}[t]\centering\footnotesize
\caption{Ablation study on SOCGraph-21. Each row turns off one component of the full agent; values are means over repeated runs ($\pm$ standard deviation when $n>1$). $\Delta$Reward is relative to the full configuration on the same tasks.}
\label{tab:ablation}
\begin{tabular}{@{}l r rrrrrrr r@{}}\toprule
Configuration & $n$ & Quality & Efficiency & Reward & Verdict acc. & FPR & FNR & Calls/task & $\Delta$Reward\\\midrule
full & 3 & 0.836\,$\pm$\,0.013 & 0.626\,$\pm$\,0.019 & 0.775\,$\pm$\,0.012 & 0.968\,$\pm$\,0.027 & 0.000\,$\pm$\,0.000 & 0.044\,$\pm$\,0.039 & 8.22\,$\pm$\,0.22 & --\\
no cache & 3 & 0.866\,$\pm$\,0.005 & 0.610\,$\pm$\,0.007 & 0.799\,$\pm$\,0.005 & 1.000\,$\pm$\,0.000 & 0.000\,$\pm$\,0.000 & 0.000\,$\pm$\,0.000 & 8.43\,$\pm$\,0.08 & +0.023\\
no validation & 2 & 0.872\,$\pm$\,0.003 & 0.618\,$\pm$\,0.001 & 0.805\,$\pm$\,0.002 & 1.000\,$\pm$\,0.000 & 0.000\,$\pm$\,0.000 & 0.000\,$\pm$\,0.000 & 8.45\,$\pm$\,0.03 & +0.029\\
no grounding & 3 & 0.843\,$\pm$\,0.027 & 0.610\,$\pm$\,0.007 & 0.779\,$\pm$\,0.023 & 0.968\,$\pm$\,0.027 & 0.000\,$\pm$\,0.000 & 0.044\,$\pm$\,0.039 & 8.49\,$\pm$\,0.19 & +0.004\\
no budget & 3 & 0.871\,$\pm$\,0.007 & 0.608\,$\pm$\,0.008 & 0.802\,$\pm$\,0.006 & 1.000\,$\pm$\,0.000 & 0.000\,$\pm$\,0.000 & 0.000\,$\pm$\,0.000 & 8.75\,$\pm$\,0.29 & +0.027\\
no method & 3 & 0.830\,$\pm$\,0.026 & 0.432\,$\pm$\,0.009 & 0.737\,$\pm$\,0.022 & 0.937\,$\pm$\,0.027 & 0.000\,$\pm$\,0.000 & 0.089\,$\pm$\,0.038 & 11.70\,$\pm$\,0.28 & -0.038\\
\bottomrule\end{tabular}\end{table*}

\begin{table*}[t]\centering\footnotesize
\caption{Paired effect of each ablation against the full agent: per-task difference, averaged over the 2 repeats with per-task results, with a bootstrap 95\% interval over the tasks (10{,}000 resamples). The third column leaves out the task whose verdict was unstable. \emph{More calls}: tasks on which the ablation used more calls than the full agent, of those where the number differed.}
\label{tab:effects}
\begin{tabular}{@{}lrrrr@{}}\toprule
Configuration & $\Delta$Reward [95\% CI] & excl.\ \texttt{linux-service-stopped} & $\Delta$Calls [95\% CI] & More calls\\\midrule
no cache & +0.017 [-0.016, +0.064] & -0.004 [-0.020, +0.009] & +0.17 [-0.29, +0.62] & 9/15\\
no validation & +0.024 [-0.004, +0.070] & +0.004 [-0.006, +0.015] & +0.17 [-0.14, +0.50] & 9/15\\
no grounding & +0.007 [-0.003, +0.018] & +0.005 [-0.005, +0.016] & +0.19 [-0.31, +0.76] & 6/12\\
no budget & +0.022 [-0.008, +0.068] & +0.001 [-0.011, +0.012] & +0.31 [-0.31, +1.02] & 9/15\\
no method & -0.031 [-0.078, +0.008] & -0.013 [-0.044, +0.015] & +3.57 [+2.43, +4.79] & 20/20\\
\bottomrule\end{tabular}\end{table*}

% Ablation: quality and efficiency per configuration (pgfplots, reads data/runs.csv).
% Drawn only when the study results exist (\ResStudyAvailable = 1).
\pgfplotstableread[col sep=comma]{data/runs.csv}\ablationruns
\begin{figure}[t]
\centering
\begin{tikzpicture}
\begin{axis}[
  ybar, bar width=6pt, width=0.95\columnwidth, height=5.4cm,
  ymin=0, ymax=1.05, ylabel={Mean over tasks}, ymajorgrids, grid style={gray!25},
  xticklabels from table={\ablationruns}{configuration}, xtick=data,
  xticklabel style={font=\scriptsize, rotate=25, anchor=north east},
  tick label style={font=\scriptsize}, legend style={font=\scriptsize, at={(0.5,1.02)}, anchor=south, legend columns=2, draw=none},
]
\addplot[fill=cA, draw=none, error bars/.cd, y dir=both, y explicit]
  table[x expr=\coordindex, y=mean_quality, y error=mean_quality_sd] {\ablationruns};
\addplot[fill=cB, draw=none, error bars/.cd, y dir=both, y explicit]
  table[x expr=\coordindex, y=mean_efficiency, y error=mean_efficiency_sd] {\ablationruns};
\legend{Quality, Efficiency}
\end{axis}
\end{tikzpicture}
\caption{Mean quality and efficiency per configuration (error bars: standard deviation over repeats).}
\label{fig:ablation}
\end{figure}


\paragraph{Effectiveness (RQ1)} The full agent reached reward \ResFullReward{}\,$\pm$\,\ResFullRewardSd{},
quality \ResFullQuality{} and verdict accuracy \ResFullAccuracy{}. No configuration judged a benign task
malicious in any run, so the false-positive rate was zero throughout. The full agent's verdict errors
(\ResFullVerdictErrors{}) were both on \texttt{linux-service-stopped}, the malicious member of the twin pair.
Across all configurations this task was cleared in 7 of the 17 runs, while its benign twin, with the same
alert, was judged correctly in all 17. The traces show why: in each erroneous run the agent observed that the
operator ran \texttt{status} and \texttt{restart} on the service seconds after stopping it, and concluded that
nothing was kept down. That sequence is the cleanup step of the recorded test, so the benign reading is
defensible from the telemetry alone; the label follows the recording's intent. The only other error was one
malicious Windows task (\texttt{firewall-rule-added}) cleared without the method prompt.

\paragraph{Platforms and verdict classes (RQ4)} Quality was \ResFullWinQuality{} on Windows tasks and
\ResFullLinQuality{} on Linux tasks. The traces show two systematic patterns that cost quality on both
platforms, more on Linux. (i) \emph{Hosts named by hostname.} Of the \TrAllMalSubs{} correct malicious
submissions, \TrAllHostByName{} named the host to isolate by its hostname instead of its entity id, which
the scorer counts as a missed core entity and an unrelated citation; only \TrAllHostById{} used the id.
(ii) \emph{Lineage relations and the session root.} On the benign Linux tasks the agent cited the
process-creation relations along the trigger's lineage, which v1.1 does not count as neutral, and stopped at
the script that explains the activity (the MOTD script, the package's post-installation script) instead of
citing the session root above it (\texttt{sshd}, \texttt{dpkg}). Hence the constant $Q=0.708$ on
\texttt{linux-dpkg-service-start} and \texttt{linux-motd-discovery}: two of three anchors and two relation
citations outside $\mathcal L\cup N$.

\begin{table}[t]\centering\footnotesize
\caption{Scoring sensitivity: mean quality of the full agent's 63 submissions rescored under alternative rules (no new model calls).}
\label{tab:sensitivity}
\resizebox{\columnwidth}{!}{%
\begin{tabular}{@{}lrrrrr@{}}\toprule
Scoring rule & All & Windows & Linux & Malicious & Benign\\\midrule
Official (v1.1) & 0.836 & 0.930 & 0.733 & 0.837 & 0.836\\
+ lineage relations neutral & 0.855 & 0.934 & 0.769 & 0.837 & 0.903\\
+ hostnames accepted & 0.876 & 0.973 & 0.769 & 0.892 & 0.836\\
+ both & 0.895 & 0.977 & 0.806 & 0.892 & 0.903\\
\bottomrule\end{tabular}}\end{table}


Table~\ref{tab:sensitivity} rescores the full agent's \TrFullSubmissions{} submissions under alternative rules,
without new model calls. Accepting hostnames raises Windows quality from \SensOfficialWindows{} to
\SensHostnamesWindows{}; treating lineage relations as neutral raises benign quality from \SensOfficialBenign{}
to \SensRelationsBenign{}; with both, overall quality rises from \SensOfficialAll{} to \SensBothAll{} and Linux
from \SensOfficialLinux{} to \SensBothLinux{}. We keep the official scorer, which is fixed and released; the
table shows how much of the measured gap reflects identifier discipline and evidence style rather than wrong
conclusions. The Linux gap that remains comes from the twin errors and from tasks where the agent drew the
operator session differently from the ground truth: on \texttt{linux-sudoers-nopasswd} it treated the script
that wrote the rule, not the login shell above it, as the session root, and named \texttt{root} rather than the
operator account.

\paragraph{Efficiency and where reward is lost (RQ2)} The full agent used \ResFullCalls{} calls per task
against a mean target $n^\star+1$ of 4.6, and exceeded the target in 58 of 63 task runs. The reward gap
decomposes into a quality part, $1-\overline{Q}=0.164$, and an efficiency part,
$\overline{Q}-\overline{R}=0.061$; on the full benchmark, unlike the Windows-only pilot, quality is the larger
loss. A task took \TrFullModelCalls{} model calls on average, with about \TrFullOutTok{} output tokens and
\TrFullCacheReadTok{} prompt tokens read from the cache; uncached input was negligible because the system
prompt and tool definitions were cached. The hard cap of 30 calls was never reached (at most 22 in any run).

\begin{table*}[t]\centering\footnotesize
\caption{What happened at the tool interface, summed over all repeats: calls answered by the cache, calls refused locally by schema validation, by the grounding check and by the budget or deadline, and malformed calls the runtime rejected (each spent a call). Tokens and model calls are means per task; \emph{Host by id/name}: correct malicious submissions naming the host by its entity id, or by its hostname.}
\label{tab:safeguards}
\begin{tabular}{@{}lrrrrrrrrrr@{}}\toprule
Configuration & Tasks & Cache & Valid. & Ground. & Budget & Runtime rej. & Output tok. & Cache-read tok. & Model calls & Host by id/name\\\midrule
full & 63 & 0 & 5 & 1 & 0 & 0 & 1{,}550 & 35{,}544 & 5.79 & 4/29\\
no cache & 63 & 0 & 4 & 1 & 0 & 0 & 1{,}547 & 38{,}281 & 5.86 & 7/28\\
no validation & 42 & 0 & 0 & 0 & 0 & 5 & 1{,}556 & 37{,}808 & 5.81 & 6/20\\
no grounding & 63 & 0 & 8 & 0 & 0 & 0 & 1{,}576 & 37{,}146 & 5.73 & 7/28\\
no budget & 63 & 1 & 7 & 0 & 0 & 0 & 1{,}627 & 40{,}162 & 6.11 & 8/29\\
no method & 63 & 0 & 12 & 0 & 0 & 0 & 2{,}024 & 32{,}241 & 6.57 & 16/10\\
\bottomrule\end{tabular}\end{table*}


\paragraph{Ablation (RQ3): the method prompt} Replacing the investigation-method prompt with the minimal prompt
is the only ablation with a consistent effect. The agent then used \ResNoMethodPairedCalls{} calls per task more
(95\% CI \ResNoMethodPairedCallsCI{}), more on \ResNoMethodMoreCalls{} of the \ResNoMethodCallsDiffer{} tasks
where the count differed (sign test \ResNoMethodCallsSignP{}), and efficiency fell from \ResFullEfficiency{} to
\ResNoMethodEfficiency{}. The paired change in mean reward was \ResNoMethodPairedReward{} (95\% CI
\ResNoMethodPairedRewardCI{}).
Quality did not fall (\ResNoMethodQuality{} against \ResFullQuality{}), for a reason the traces make visible:
the method prompt tells the agent not to look up the host named in the alert, so the full agent rarely held the
host's entity id (\TrFullHostById{} of \TrFullMalSubs{} correct malicious submissions used it, against
\TrNoMethodHostById{} of \TrNoMethodMalSubs{} without the method). An efficiency instruction thus cost quality
through the interface, offsetting what the method gains elsewhere. Without the method the agent also sent more
malformed calls (\TrNoMethodValidation{} refused by validation, against \TrFullValidation{}) and used about
30\% more output tokens.

\paragraph{Ablation (RQ3): the local safeguards} Switching off the cache, validation, grounding check or budget
produced no measurable change. Three of these ablations score about 0.02--0.03 above the full agent in mean
reward, but the difference is the twin task: the full agent cleared \texttt{linux-service-stopped} in two of
three repeats, and each miss removes about 0.04 from the mean over 21 tasks. Leaving that task out, every
paired difference lies within $\pm0.004$ and every interval contains zero (Table~\ref{tab:effects}). The
per-configuration error counts on that task are consistent with one common rate: with the pooled 7 errors in 17
runs, two or more errors in three runs has probability of about 0.37 and none about 0.20.

Table~\ref{tab:safeguards} shows why the safeguards had nothing to win. In \TrFullTasks{} full-agent tasks the
cache answered \TrFullCache{} calls, validation refused \TrFullValidation{} malformed calls (a search query
shorter than the schema's minimum, a submission serialised as an invalid JSON string), the grounding check
fired \TrFullGrounding{} time (an invented host identifier, \texttt{host-workstation5}, which the agent then
dropped) and the budget never refused a call. Without validation, the same kinds of malformed call reached the
runtime \TrNoValidationRuntimeSchema{} times in 42 tasks and were rejected there, each costing a call: about
0.1 calls per task, too few to move efficiency. The grounding check also missed the most frequent identifier
error: it accepts any string that occurred in the evidence, and hostnames do occur there. A check that accepts
only values seen in identifier fields would have flagged all \TrAllHostByName{} hostname citations.
\else
The full study runs every configuration of \S\ref{sec:experiments} three times on all 21 tasks. Its results
are generated into Tables~\ref{tab:ablation}--\ref{tab:pertask} and a comparison figure by the release
script from the run folders, and this subsection will report them.
\fi

\ifnum\ResStudyAvailable=0
\begin{table*}[t]\centering\footnotesize
\caption{Ablation study on SOCGraph-21. Each row turns off one component of the full agent; values are means over repeated runs ($\pm$ standard deviation when $n>1$). $\Delta$Reward is relative to the full configuration on the same tasks.}
\label{tab:ablation}
\begin{tabular}{@{}l r rrrrrrr r@{}}\toprule
Configuration & $n$ & Quality & Efficiency & Reward & Verdict acc. & FPR & FNR & Calls/task & $\Delta$Reward\\\midrule
full & 3 & 0.836\,$\pm$\,0.013 & 0.626\,$\pm$\,0.019 & 0.775\,$\pm$\,0.012 & 0.968\,$\pm$\,0.027 & 0.000\,$\pm$\,0.000 & 0.044\,$\pm$\,0.039 & 8.22\,$\pm$\,0.22 & --\\
no cache & 3 & 0.866\,$\pm$\,0.005 & 0.610\,$\pm$\,0.007 & 0.799\,$\pm$\,0.005 & 1.000\,$\pm$\,0.000 & 0.000\,$\pm$\,0.000 & 0.000\,$\pm$\,0.000 & 8.43\,$\pm$\,0.08 & +0.023\\
no validation & 2 & 0.872\,$\pm$\,0.003 & 0.618\,$\pm$\,0.001 & 0.805\,$\pm$\,0.002 & 1.000\,$\pm$\,0.000 & 0.000\,$\pm$\,0.000 & 0.000\,$\pm$\,0.000 & 8.45\,$\pm$\,0.03 & +0.029\\
no grounding & 3 & 0.843\,$\pm$\,0.027 & 0.610\,$\pm$\,0.007 & 0.779\,$\pm$\,0.023 & 0.968\,$\pm$\,0.027 & 0.000\,$\pm$\,0.000 & 0.044\,$\pm$\,0.039 & 8.49\,$\pm$\,0.19 & +0.004\\
no budget & 3 & 0.871\,$\pm$\,0.007 & 0.608\,$\pm$\,0.008 & 0.802\,$\pm$\,0.006 & 1.000\,$\pm$\,0.000 & 0.000\,$\pm$\,0.000 & 0.000\,$\pm$\,0.000 & 8.75\,$\pm$\,0.29 & +0.027\\
no method & 3 & 0.830\,$\pm$\,0.026 & 0.432\,$\pm$\,0.009 & 0.737\,$\pm$\,0.022 & 0.937\,$\pm$\,0.027 & 0.000\,$\pm$\,0.000 & 0.089\,$\pm$\,0.038 & 11.70\,$\pm$\,0.28 & -0.038\\
\bottomrule\end{tabular}\end{table*}

\begin{table*}[t]\centering\footnotesize
\caption{Paired effect of each ablation against the full agent: per-task difference, averaged over the 2 repeats with per-task results, with a bootstrap 95\% interval over the tasks (10{,}000 resamples). The third column leaves out the task whose verdict was unstable. \emph{More calls}: tasks on which the ablation used more calls than the full agent, of those where the number differed.}
\label{tab:effects}
\begin{tabular}{@{}lrrrr@{}}\toprule
Configuration & $\Delta$Reward [95\% CI] & excl.\ \texttt{linux-service-stopped} & $\Delta$Calls [95\% CI] & More calls\\\midrule
no cache & +0.017 [-0.016, +0.064] & -0.004 [-0.020, +0.009] & +0.17 [-0.29, +0.62] & 9/15\\
no validation & +0.024 [-0.004, +0.070] & +0.004 [-0.006, +0.015] & +0.17 [-0.14, +0.50] & 9/15\\
no grounding & +0.007 [-0.003, +0.018] & +0.005 [-0.005, +0.016] & +0.19 [-0.31, +0.76] & 6/12\\
no budget & +0.022 [-0.008, +0.068] & +0.001 [-0.011, +0.012] & +0.31 [-0.31, +1.02] & 9/15\\
no method & -0.031 [-0.078, +0.008] & -0.013 [-0.044, +0.015] & +3.57 [+2.43, +4.79] & 20/20\\
\bottomrule\end{tabular}\end{table*}

\fi
\begin{table*}[t]\centering\footnotesize
\caption{Quality and verdict accuracy by platform and by expected verdict, over the repeats with per-task results (mean $\pm$ standard deviation across repeats).}
\label{tab:breakdown}
\begin{tabular}{@{}l rrrrrrrr@{}}\toprule
Configuration & \multicolumn{2}{c}{Windows} & \multicolumn{2}{c}{Linux} & \multicolumn{2}{c}{Malicious} & \multicolumn{2}{c}{Benign}\\
\cmidrule(lr){2-3}\cmidrule(lr){4-5}\cmidrule(lr){6-7}\cmidrule(lr){8-9}
 & Qual. & Acc. & Qual. & Acc. & Qual. & Acc. & Qual. & Acc.\\\midrule
full & 0.930\,$\pm$\,0.017 & 1.00\,$\pm$\,0.00 & 0.733\,$\pm$\,0.034 & 0.93\,$\pm$\,0.06 & 0.837\,$\pm$\,0.020 & 0.96\,$\pm$\,0.04 & 0.836\,$\pm$\,0.004 & 1.00\,$\pm$\,0.00\\
no cache & 0.938\,$\pm$\,0.012 & 1.00\,$\pm$\,0.00 & 0.787\,$\pm$\,0.010 & 1.00\,$\pm$\,0.00 & 0.879\,$\pm$\,0.008 & 1.00\,$\pm$\,0.00 & 0.833\,$\pm$\,0.004 & 1.00\,$\pm$\,0.00\\
no validation & 0.941\,$\pm$\,0.009 & 1.00\,$\pm$\,0.00 & 0.795\,$\pm$\,0.005 & 1.00\,$\pm$\,0.00 & 0.885\,$\pm$\,0.002 & 1.00\,$\pm$\,0.00 & 0.838\,$\pm$\,0.003 & 1.00\,$\pm$\,0.00\\
no grounding & 0.942\,$\pm$\,0.005 & 1.00\,$\pm$\,0.00 & 0.735\,$\pm$\,0.052 & 0.93\,$\pm$\,0.06 & 0.845\,$\pm$\,0.036 & 0.96\,$\pm$\,0.04 & 0.840\,$\pm$\,0.007 & 1.00\,$\pm$\,0.00\\
no budget & 0.932\,$\pm$\,0.007 & 1.00\,$\pm$\,0.00 & 0.803\,$\pm$\,0.010 & 1.00\,$\pm$\,0.00 & 0.887\,$\pm$\,0.007 & 1.00\,$\pm$\,0.00 & 0.831\,$\pm$\,0.009 & 1.00\,$\pm$\,0.00\\
no method & 0.920\,$\pm$\,0.042 & 0.97\,$\pm$\,0.05 & 0.731\,$\pm$\,0.017 & 0.90\,$\pm$\,0.00 & 0.826\,$\pm$\,0.038 & 0.91\,$\pm$\,0.04 & 0.839\,$\pm$\,0.002 & 1.00\,$\pm$\,0.00\\
\bottomrule\end{tabular}\end{table*}

\begin{table*}[t]\centering\scriptsize
\caption{Mean reward per task and configuration, over the repeats with per-task results. A dagger marks a configuration whose verdict was wrong in at least one repeat.}
\label{tab:pertask}
\begin{tabular}{@{}lllrrrrrr@{}}\toprule
Task & Platform & Expected & full & no cache & no validation & no grounding & no budget & no method\\\midrule
\texttt{bits-download} & windows & malicious & 0.857 & 0.866 & 0.835 & 0.857 & 0.819 & 0.748\\
\texttt{firewall-rule-added} & windows & malicious & 0.792 & 0.812 & 0.778 & 0.789 & 0.777 & 0.569$^\dagger$\\
\texttt{gpsvc-service-host} & windows & benign & 0.779 & 0.796 & 0.784 & 0.792 & 0.821 & 0.710\\
\texttt{hta-startup-folder} & windows & malicious & 0.860 & 0.847 & 0.843 & 0.831 & 0.827 & 0.834\\
\texttt{linux-account-created} & linux & malicious & 0.801 & 0.800 & 0.779 & 0.771 & 0.760 & 0.762\\
\texttt{linux-dpkg-service-start} & linux & benign & 0.661 & 0.661 & 0.661 & 0.661 & 0.661 & 0.776\\
\texttt{linux-kernel-module} & linux & malicious & 0.817 & 0.797 & 0.823 & 0.809 & 0.811 & 0.827\\
\texttt{linux-ld-preload-hijack} & linux & malicious & 0.691 & 0.674 & 0.656 & 0.712 & 0.714 & 0.705\\
\texttt{linux-motd-discovery} & linux & benign & 0.661 & 0.661 & 0.661 & 0.661 & 0.657 & 0.648\\
\texttt{linux-motd-tmp-cleanup} & linux & benign & 0.703 & 0.727 & 0.729 & 0.727 & 0.699 & 0.617\\
\texttt{linux-service-stopped} & linux & malicious & 0.252$^\dagger$ & 0.830 & 0.811 & 0.276$^\dagger$ & 0.830 & 0.000$^\dagger$\\
\texttt{linux-shadow-read} & linux & malicious & 0.772 & 0.767 & 0.763 & 0.772 & 0.799 & 0.772\\
\texttt{linux-sudoers-nopasswd} & linux & malicious & 0.588 & 0.586 & 0.644 & 0.585 & 0.597 & 0.613\\
\texttt{linux-wiper-shred} & linux & malicious & 0.825 & 0.720 & 0.812 & 0.772 & 0.818 & 0.717\\
\texttt{logon-script-persistence} & windows & malicious & 0.876 & 0.864 & 0.882 & 0.883 & 0.900 & 0.916\\
\texttt{lsass-dump-comsvcs} & windows & malicious & 0.904 & 0.937 & 0.934 & 0.931 & 0.921 & 0.900\\
\texttt{lsass-query-by-svchost} & windows & benign & 1.000 & 1.000 & 1.000 & 1.000 & 1.000 & 0.950\\
\texttt{python-http-server} & windows & malicious & 0.777 & 0.815 & 0.838 & 0.833 & 0.794 & 0.844\\
\texttt{run-key-and-beacon} & windows & malicious & 0.876 & 0.885 & 0.866 & 0.907 & 0.900 & 0.897\\
\texttt{service-binpath-hijack} & windows & malicious & 0.837 & 0.837 & 0.849 & 0.838 & 0.851 & 0.820\\
\texttt{w32time-runtime-key} & windows & benign & 0.950 & 0.892 & 0.950 & 0.950 & 0.892 & 0.856\\
\bottomrule\end{tabular}\end{table*}


\subsection{Shortcuts and evasion (RQ5)}
\label{sec:cues}
The study's traces show that the agent does read label cues. Of the \TrAllSubmissions{} submissions,
\TrAllSourceCue{} (\TrAllSourceCuePct\%) justify the verdict partly by the source recording, calling the
activity a simulation or naming the attack dataset from the alert's source field. Most of these are correct
malicious verdicts (\TrAllSourceCueRight{}). The other \TrAllSourceCueWrong{} are every wrong verdict of the
study (the seven cleared \texttt{linux-service-stopped} runs and the cleared \texttt{firewall-rule-added} run).
In these the agent named the source only as a caveat (``the alert source is a Splunk attack-simulation
dataset, so the label is a judgement call'') and decided from the graph, where it saw no lasting effect. The
cue therefore did not override the evidence when the two disagreed; whether it helped where they agreed is what
the cue-controlled variants test.
\ifnum\ResCueAvailable=1
\begin{table*}[t]\centering\footnotesize\setlength{\tabcolsep}{4pt}
\caption{Shortcuts and evasion: the full agent on \bench v1.1 and on its cue-controlled variants (same graphs, identifiers and answers). Means $\pm$ standard deviation over repeats; $\Delta$Reward is the paired per-task difference to v1.1 with a bootstrap 95\% interval over tasks; \emph{Missed attacks}: malicious task runs judged benign.}
\label{tab:cues}
\begin{tabular}{@{}lrrrrrrrr@{}}\toprule
Condition & $n$ & Quality & Efficiency & Reward & Verdict acc. & Calls/task & $\Delta$Reward [95\% CI] & Missed attacks\\\midrule
SOCGraph-21 v1.1 & 3 & 0.836\,$\pm$\,0.013 & 0.626\,$\pm$\,0.019 & 0.775\,$\pm$\,0.012 & 0.968\,$\pm$\,0.027 & 8.22\,$\pm$\,0.22 & -- & 2/45\\
\emph{blind} & 3 & 0.836\,$\pm$\,0.006 & 0.586\,$\pm$\,0.013 & 0.769\,$\pm$\,0.007 & 0.952\,$\pm$\,0.000 & 9.00\,$\pm$\,0.21 & -0.006 [-0.037, +0.015] & 3/45\\
\emph{mimicry} & 3 & 0.754\,$\pm$\,0.044 & 0.572\,$\pm$\,0.005 & 0.694\,$\pm$\,0.038 & 0.857\,$\pm$\,0.048 & 9.22\,$\pm$\,0.15 & -0.081 [-0.179, -0.007] & 9/45\\
\bottomrule\end{tabular}\end{table*}

% Interpretation of the variant results goes here.
\else
The runs on the \emph{blind} and \emph{mimicry} variants (\S\ref{sec:variants}) are specified in
\S\ref{sec:experiments}; their results will be added to this subsection.
\fi

\subsection{Engineering observations}
\label{sec:engineering}

\paragraph{Policy-compatible rationale fields} With the runtime's field names exposed as they were, every task
of an initial mock run was declined at its first request with the refusal category
\texttt{reasoning\_extraction}. After presenting the fields as \texttt{purpose} and \texttt{summary} and removing the instruction to
state hypotheses from the first message (\S\ref{sec:alias}), all tasks completed, and the runtime still
received its required \texttt{reasoning} fields.

\paragraph{Parallel execution} On a simulated runtime that allowed two tasks in flight, the agent was started
with four workers. It learned the limit from the first refused leases (the log shows the limit lowered first
to three, then to two, as refusals arrived), never ran more than two tasks at once afterwards, and completed
all 16 tasks of that run in 112\,s with the real model, without a failed task.

\paragraph{Benign twins} The two Linux service tasks share an identical alert (``Service state changed by
systemctl as root'') and graph, and differ only in the lineage of the triggering process: an operator's login
shell through \texttt{sudo} versus the package manager's post-installation script. They test whether an agent
reads provenance rather than the alert text. The agent never called the package manager's activity malicious,
but it cleared the operator's service stop in 7 of 17 runs: an interactive administrator stopping a web server
is an ordinary event in isolation, and only the lab context makes it an attack. Twin pairs thus expose the
direction in which an agent errs when provenance alone is ambiguous.
